\mainmatter
\silentchapter{Proposisjonar}{proposisjonar}

\prop{1}{}{Formverda er alt som er tilfelle.}

\prop{1.1}{o}{Formverda er totaliteten av realiserte konfigurasjonar, ikkje ting. Alt som har form, har eksakt denne forma. At ein konfigurasjon tek éi form framfor andre moglege, krev ei ekstern årsak — ei årsak som kviler i relasjonen mellom den manifesterte konfigurasjonen og dei urealiserte moglegheitene. Dette er ein relasjonell, ikkje absolutt, eigenskap.}

\prop{1.2}{d}{Eit \textbf{objekt} er den enklaste bestanddelen i ei form. Objektet er udeleleg og uforanderleg; det utgjer substansen i formverda. Alle moglege former er alt gjevne med objekta — moglegheita eksisterer logisk fordi objekta finst, uavhengig av tenking.}

\prop{1.21}{d}{Ein \textbf{konfigurasjon} er ei spesifikk samansetjing av objekt. Forma er strukturen til konfigurasjonen.}

\prop{1.3}{d}{\textbf{Formrommet} (\textit{morphospace}) til ein klasse er mengda av alle hennar moglege konfigurasjonar. Kvart realisert objekt okkuperer eitt eksakt punkt i dette n-dimensjonale rommet.}

\prop{1.31}{d}{Det empiriske formrommet er alltid ein projeksjon. Inga endeleg parametermengd fangar den latente kompleksiteten fullt ut. Projeksjonen vel: ho lyser opp somme relasjonar og mørklegg andre. Det finst ingen nøytral projeksjon. Projeksjonen gjer empirien mogleg, men konstituerer ikkje forma — forma eksisterer uavhengig av målinga.}

\prop{1.32}{d}{Klassen avgrensar formrommet funksjonelt. Klassen set grensene; seleksjonstrykket fordelar innhaldet.}

\prop{1.33}{d}{Morforommet er ein algebra lukka under union, differanse og euklidsk transformasjon. Algebraen definerer horisonten for alle generative operasjonar.}

\prop{1.4}{d}{Formrommet deler seg topologisk i tre regionar: busette, opne og forbodne.}

\prop{1.41}{o}{Dei busette regionane utgjer det historiske arkivet. Dei fungerer som ankerpunkt for framtidig navigasjon; tyngda til ankerpunktet veks med okkupasjonstid og konvergerande tradisjonar. Uavhengig konvergens provar eit reelt optimum sterkare enn éin einskild tradisjon.}

\prop{1.42}{o}{Dei opne regionane utgjer det \textit{tilstøytande moglege}. Dei er tilgjengelege, men uaktualiserte. Det tilstøytande moglege er endeleg — spent ut mellom busette naboar og forbodne barrierar. Det busette tiltrekkjer; det forbodne fråstøyter. Det realiserte rommet oppstår i spennet mellom dei to.}

\prop{1.43}{d}{Dei forbodne regionane er utelukka av materielle, strukturelle eller energetiske restriksjonar. Forbodet følgjer vilkåra, ikkje regionen — grensa mellom det opne og det forbodne er dynamisk. Eit materialskifte kan opne ein forboden region. Grensa tilhøyrer vilkåra, ikkje forma.}

\prop{1.44}{d}{Dei tre regionane utgjer ein epistemisk partisjon: det kjende (K), det tenkelege ukjende (C\msym{∖}K), og det utenkjelege. Grensene mellom sonene er epistemiske, ikkje ontologiske: eit skifte i vilkåra omkalfatrar partisjonen.}

\silentchapter{2 Ikkje alle posisjonar er like sannsynlege}{ikkje alle posisjonar er like sannsynlege}

\prop{2}{}{Ikkje alle posisjonar er like sannsynlege.}

\prop{2.1}{d}{Eit \textbf{seleksjonstrykk} endrar sannsynsfordelinga i formrommet. Trykket er inga regel, men ein gradient som aukar sannsynet for visse posisjonar. Eitt seleksjonstrykk avgjer aldri posisjonen åleine — det avgrensar kvar posisjonen ikkje kan vere. Forma står att når trykka har utelukka alt anna. Forma er ein rest.}

\prop{2.11}{t}{Eit trykk som inkje utelukkar, er inkje trykk. Eit trykk som utelukkar alt unnateke éin region, er determinisme. Observerte trykk ligg mellom desse ytterpunkta.}

\prop{2.12}{a}{Fleire seleksjonstrykk verkar samstundes i kvar funksjonell klasse. Einsidige formforklaringar er difor ufullstendige: med berre eitt trykk ville alle former vore identiske eller tilfeldig fordelte langs éin akse. Ikkje-uniformiteten krev fleire uavhengige trykk.}

\prop{2.121}{t}{Eit latent trykk verkar framleis. Uvarierte trykk lèt seg knapt telje; metoden krev variasjon i vilkåra.}

\prop{2.13}{a}{Minst to trykk peikar i motstridande retningar i minst éin region. Trykk i motstridande retningar utelukkar globale optimum. Kvar realisert form er eit kompromiss — den einaste moglege balansen under det gjevne trykksettet.}

\prop{2.131}{t}{Formskilnader syner at kreftene lèt seg balansere ulikt. Omgrepet «suboptimal» føreset ein referanse; utan eit spesifikt trykksett er det ikkje falsifiserbart.}

\prop{2.14}{t}{Funksjonen definerer klassen og formrommet. Funksjonen er konstant innanfor klassen og ber difor inga informasjon om kvifor éi form oppstod. Dei resterande seleksjonstrykka driv all formvariasjon i klassen.}

\prop{2.141}{i}{Stolen og sofaen utgjer ulike klasser med ulik funksjon. To stolar med ulik stil har same funksjon. At stilen forklarer meir formvariasjon enn materialet, falsifiserer den funksjonalistiske doktrinen.}

\prop{2.2}{d}{\textbf{Materialaffordansen} er operasjonane materialet tillèt og hindrar. Signaturen bind sannsynsfordelinga, ikkje einskildobjektet. Signaturen er latent — repertoaret slår ikkje tvingande ut i geometri.}

\prop{2.21}{o}{Materialet kan finnast utan algebraen; operasjonane endrar seg, ikkje materialet. Stål fanst i hundreår før røyrstolen fordi operasjonen mangla. Formhistoria er systematisk langsamare enn materialhistoria. Nye materiale arvar det gamle formspråket inntil eigne affordansar pressar dei ut.}

\prop{2.22}{t}{Ein samlevariabel predikerer formvariasjon betre enn einskildtrykk. Stilperioden er ein slik samlevariabel — men stilen er inga forklaring, berre ein proxy som absorberer alle trykk. Stilkategorien skildrar stansen, ikkje årsaka. Forveksling av deskriptiv og kausal kraft forfalskar formhistoria. Stilperioden er eit symptom på underliggjande seleksjonstrykk.}

\silentchapter{3 Seleksjonstrykka produserer eit landskap}{seleksjonstrykka produserer eit landskap}

\prop{3}{}{Seleksjonstrykka spenner ut eit landskap over formrommet.}

\prop{3.1}{d}{\textbf{Tilpassingslandskapet} er vektorsummen av alle verksame seleksjonstrykk i formrommet. Kvar posisjon held ein verdi: graden av samtidig forløysing for alle trykk. Kvart trykk kastar ei vekt som fluktuerer i tid; vektene er bundne.}

\prop{3.11}{t}{Landskapet lèt seg berre lodde retrospektivt. For ein gjeven epoke frys vektene ut frå formene i kjernen. Landskapet yter lokal prediksjon — det peikar ut eignene som trekkjer til seg former — men ingen einskildform er garantert. Prediksjonshorisonten bøyer seg etter fluktuasjonen i landskapet: i stasen er siktlinja uendeleg; i bifurkasjonen brotnar ho.}

\prop{3.12}{t}{Av 2.12, 2.13 og 3.1 følgjer: tilpassingsfunksjonen har talrike lokale maksima. Landskapet er eit taggete massiv av mange haugar. Ein haug er staden der kvart avsteg krev eit fall i tilpassing; dalar er tilstandane formene skyr — dalane er opne, men bebuarane deira går under.}

\prop{3.121}{t}{Landskapet med éin universell topp er eit teoretisk unntak. Empirien knuser funksjonalismens postulat om den eine, perfekte forma. Den multimodale fordelinga og avstanden mellom klyngene teiknar konturane av det faktiske tilpassingslandskapet.}

\prop{3.2}{d}{Stabiliteten til ein haug svarar til brattleiken på dalveggane. Bratte veggar tvingar fram kanalisering: forma støyter vekk perturberingar. Nokre seleksjonstrykk regjerer over andre — dominante trykk faldar rommet til tronge korridorar. Trykkhierarkiet framkallar eit kanalhierarki.}

\prop{3.21}{o}{Djupe kanalar frys dimensjonane; grunne kanalar rommar fluktuasjonen. Asymmetrien mellom frosen ryggrad og fluktuerande ornament avtrykkjer landskapet sin eigen natur.}

\prop{3.3}{d}{Ein \textbf{stil} er forma si sverming kring eit felles tyngdepunkt. Sidan stilen er ein sverm og ikkje ein kategori, er grensene naudsynleg uskarpe — overgangen mellom to klynger i eit kontinuerleg rom er alltid sløra. Stilperiodar utgjer gradientar, ikkje hermetiske øyar.}

\prop{3.31}{t}{Kjernen i ein epoke slektar meir på framtida enn på si eiga randsone. Sidan stilen oppstår statistisk, herskar det alltid formell tvil ved randene. Grensene svinn; diffusheita er absolutt.}

\prop{3.4}{t}{Formgjevinga skapar aldri sitt eige landskap. Landskapet fylgjer tvingande av vilkåra. Å reise eit bygg på ein topp er ei oppdaging, ikkje eit under.}

\prop{3.41}{t}{Uavhengig konvergens stadfester at haugen spring ut av vilkåra, ikkje av utveksling. Same grunnkrav framkallar same gradient mot same spiss. Isolerte tradisjonar navigerer den same kløfta utan å utveksle eit ord. Isometriske trykk støyper same faste form uavhengig av avsendar.}

\silentchapter{4 Landskapet er dynamisk}{landskapet er dynamisk}

\prop{4}{}{Landskapet er dynamisk.}

\prop{4.1}{a}{Seleksjonstrykka kviler aldri. Av 3.1 og 4.1: Landskapet er ei uroleg flate. Haugar stig, søkk, flyttar seg, deler seg eller smeltar saman. Ein optimal posisjon vert suboptimal straks vilkåra endrar seg — optimal tilpassing er lokal i tid.}

\prop{4.11}{o}{Ein ny teknologi flyttar grensa for det forbodne; ein ny marknad endrar kva landskapet belønnar; eit nytt materiale utvidar operasjonsrommet. Kvart slikt skifte grip inn i landskapet.}

\prop{4.12}{o}{Medianen for proporsjonen høgde til breidde fell systematisk over fem hundreår. Proporsjonen 1.36 er ikkje meir ergonomisk enn 1.88; begge høver for menneskekroppen. Dette er formdrift, ikkje funksjonsdrift.}

\prop{4.2}{o}{Endringa er diskontinuerleg. Lange periodar med stase bryt saman i brå topologiske skifte. Forløpet er alltid: eksplosiv radiasjon inn i ein ny region, deretter konvergens mot nye attraktorar. Episodisk diskontinuitet bryt ikkje med den underliggande dynamikken — ho er signaturen til denne dynamikken.}

\prop{4.21}{o}{Brotet etter 1900 syner ei rekurrensmatrise der ingen periode før 1900 liknar nokon etter. Formhistoria gjentek seg aldri over dette brotet.}

\prop{4.3}{t}{Landskapet har minne. Kvar realisert form set spor som endrar seleksjonstrykka for den neste. Form og landskap endrar kvarandre.}

\prop{4.31}{t}{All formgjeving er omformgjeving; agenten startar aldri med blanke ark. Å starte frå ein posisjon er å arve alle implikasjonane hennar — dei opne og dei stengde naboregionane. Utgangspunktet er aldri nøytralt. Stiavhengigheit er eit fundamentalt trekk ved formhistoria, ingen anomali.}

\prop{4.4}{o}{Formrommet ekspanderer kumulativt. Nye regionar vert okkuperte; dei gamle vert sjeldan heilt forlatne. Det kumulative konvekse hylstervolumet veks monotont gjennom alle periodar, utan eitt einaste tilbakesteg.}

\prop{4.41}{t}{Eit fast funksjonsenvelop tvingar fram metning, men nyheitsraten fell ikkje monotont — han skyt fart att når nye materiale trer inn. Det tilstøytande moglege ekspanderer aktivt; det er inkje lager som tømmest. Materialstraumar driv landskapsendringa.}

\prop{4.42}{t}{Det tilstøytande moglege er C\msym{∖}K-grensa: mengda av former som er konseptuelt tilgjengelege men enno ukjende som realiserbare. Kvart nytt materiale opnar nokre regionar og låser andre. Materialstraumen frå 1500 til 2025 syner tydelege seleksjonsbølgjer — kvar bølgje er ei brå kohorthending, ingen gradvis overgang.}

\prop{4.5}{o}{Når eitt seleksjonstrykk dominerer, kollapsar landskapet mot éin attraktor. Eit eineveldande trykk forklarer ikkje form — det er fråværet av forklaring. Formrommets diversitet speglar mangfaldet av aktive trykk; ein monokultur provar at eitt trykk har vorte eineveldande.}

\silentchapter{5 Det finst agentar}{det finst agentar}

\prop{5}{}{Verda rommar agentar som responderer på landskapet.}

\prop{5.1}{a}{Kvar funksjonell klasse har minst eitt system som navigerer landskapet via negativ tilbakekopling.}

\prop{5.11}{d}{Ein \textbf{agent} er ein operator definert av trippelet: måltilstand, avstandsmåling, justeringsmekanisme. Definisjonen føreset tilbakekopling. Agens krev berre funksjonell organisering, verken medvit, intensjon eller nervesystem — eit system som navigerer substratblint mot eit mål, er ein agent. Ein rullande stein er ingen agent; rørsla manglar avstandsmåling. Informasjonsflyt skil agentar frå passiv materie.}

\prop{5.12}{o}{Tre målbare dimensjonar differensierer agentar: representasjonsspennvidde i formrommet, tilbakekoplingsfrekvens og regeloppdateringsdjupne. Agens er stratifisert — termostaten og arkitekten skil seg berre i navigasjonskompleksitet.}

\prop{5.13}{t}{Agenten responderer utelukkande på gradienten i landskapet, ikkje på forma. Agenten oppfinn ikkje måltilstanden; målet er ein singularitet agenten oppdagar i landskapet. Navigasjon krev inga global innsikt — agenten treng berre integrere den lokale gradienten. Agentar som persiperer same lokale gradient, induserer same trajektorie. Lokal kompetanse produserer global manifestasjon.}

\prop{5.2}{d}{\textbf{Kognitiv lyskjegle} er delmengda av formrommet agenten kan representere og manipulere. Alt utanfor lyskjegla manglar kausal eksistens. Lyskjegler spenner over vilkårlege skalaer — frå cellulær millisekundrespons til tiårslang distribuert konsensus.}

\prop{5.21}{t}{Lyskjegla dikterer oppløysinga til formrommet. Agentens diskriminering av sub-geometriar avgjer det tilgjengelege transformasjonsrommet. Agenten persiperer utelukkande affordansar, ikkje former — ein affordans er den lokale asymmetrien manifest i agentens lyskjegle. Utviding av lyskjegla aktualiserer latente affordansar; innsnevring deaktiverer reelle.}

\prop{5.22}{d}{Fem kanoniske operasjonar muterer lyskjegla: under ekspansjon innlemmar ho nye topologiske regionar; under refraksjon materialiserer sub-komponentar seg i det etablerte rommet; under translasjon forskyvst ho langs formrommets aksar; under kollaps singulariserer ho mot éin urokkeleg posisjon; under spissing bevarer ho volumet medan operasjonspresisjonen asymptoterer mot ei grense. Kausal kontakt og rekursiv justering er dei einaste vilkåra.}

\prop{5.3}{d}{Agentens essensielle funksjon er direkte manipulering av systemgeometrien. Vektorrommet av alle konfigurasjonar lukka under union og snitt konstituerer ein \textbf{formalgebra}; \textbf{formgrammatikken} utfører diskrete overgangar i denne geometrien. Grammatikken gjenkjenner og overskriv konfigurasjonelle undergrupper.}

\prop{5.31}{t}{Transformasjonsreglar fyrer utelukkande på strukturell innleiring — mønstergjenkjenning opererer uavhengig av opphavshistorie. Historikken er oppheva fordi reglane opererer utelukkande i notid; inndatageometrien er regelens einaste premiss, og slettar dermed produksjonsstien.}

\prop{5.32}{t}{Grammatikkoperasjonen er ein C\msym{→}C-partisjon: ho utvidar konseptet ved å leggje til eigenskapar. Realisering er ein C\msym{→}K-operasjon: ho avgjer konseptet.}

\prop{5.33}{t}{Grammatikken induserer ein naboskap for kvar gyldig form — mengda av konfigurasjonar tilgjengelege gjennom éin diskret operasjon. Naboskapen utgjer formens umiddelbare kausale fridomsrom. Naboskapens asymmetri utelukkar garanterte returvegar; formrommet utgjer ein retta graf. Formsuperposisjon kan frigjere delformer uforutsette av grunnkomponentane — forma overgår operasjonane som framkalla ho.}

\prop{5.34}{t}{Grammatikken definerer moglegheitsrommet; landskapstrykket fastset den endelege trajektorien. Grammatikk og gradient kolliderer i agentens operasjonelle singularitet. Landskapet tildeler vekter til naboskapens utstaka vegar; agenten forløyser konfigurasjonen som mest radikalt reduserer lokal spenning. Lyskjeglas sensur av innleiringar gjer generativ makt til ein funksjon av observasjonsvindauget.}

\prop{5.4}{d}{Substratet fungerer som ein null-ordens agent. Substratet søkjer inga maksimering, men vetoar konfigurasjonar og gjev resonans til andre gjennom ufravikelege avgrensingar. Substratet innrissar landskapets ontologiske furer.}

\prop{5.41}{i}{Biologisk vev navigerer gjennom bioelektrisk polarisering; menneskeleg praksis navigerer via materiell friksjon; ein nevral diffusjonsmodell navigerer det latente rommet via støyfjerning; ein slimsopp finn det optimale rutenettet via oscillerande samandragingar; ein stor språkmodell navigerer det semantiske landskapet via vektoriell merksemd. Navigasjonsmekanikken er universell trass i den fysiske implementeringas vilkårlegheit. Kompetansen deira er formelt ekvivalent: han målast i evna til å minimere spenning langs dei lokale aksane i deira respektive lyskjegler.}

\prop{5.5}{t}{Navigasjonskompetanse akkumulerer seg som strukturert informasjon. Erobra posisjonar utvidar lyskjegla for framtidig persepsjon.}

\prop{5.51}{d}{Læring er den irreversible morfologiske ekspansjonen av tilgjengeleg formrom. Læring, evolusjonær drift og nettverksruting utfører isomorfe gradientnedstigingar over ulike substrat. Prosessane minskar den topologiske avstanden mellom struktur og landskap.}

\prop{5.6}{d}{\textbf{Omsorg} (\textit{care}) er mengda av aksar i formrommet agenten aktivt kan representere og justere for.}

\prop{5.61}{t}{Omsorg er intelligensens innhald; intelligens er omsorgas form.}

\prop{5.62}{t}{Ein agent utan omsorg for ein dimensjon, er blind for landskapets gradient langs denne aksen. Kvantifisert omsorg svarar til det topologiske volumet av agentens lyskjegle. Å utvide lyskjegla er ein etisk operasjon: det er å ta fleire krefter med i det formelle kompromisset.}

\silentchapter{6 Forma oppstår mellom agentane}{forma oppstår mellom agentane}

\prop{6}{}{Forma oppstår i feltet mellom agentar.}

\prop{6.1}{t}{Av 2.12 og 5.1: Fleire trykk verkar alltid. Ulike agentar responderer på ulike trykk. Kvar form er eit poly-agentisk kompromiss. Materialaffordansen, reiskapen, marknaden og kulturen utgjer uavhengige gradientar. Ingen einskild agent dikterer morfologien — ho trer fram i skjeringspunktet mellom dei.}

\prop{6.11}{t}{Kvart nivå løyser problem kompetent innanfor si eiga lyskjegle. Lågare komponentar navigerer lokale gradientar og fremjar mål dei manglar kapasitet til å representere. Reduksjonisme og holisme er komplementære skildringsnivå: begge er sanne, ingen er tilstrekkeleg åleine.}

\prop{6.12}{t}{Agens er multiskalar: frå molekylære sveisepunkt til globale distribusjonsnettverk. Kvar skala har sin eigen kompetanse og si eiga lyskjegle. Formproduksjon er distribuert C\msym{↔}K-ekspansjon utført av grammatikkoperasjonar over kompetanseskalaer.}

\prop{6.13}{t}{Landskapet avgjer kva reglar som fyrer, uavhengig av nominell kontroll. Agenten vel; landskapet filtrerer. Kontrollen er alltid delt.}

\prop{6.14}{o}{Endeleg bandbreidd avgrensar all kommunikasjon. Låg bandbreidd gjev grov, stokastisk navigasjon; høg bandbreidd tillèt eksplisitt artikulering. Bandbreidd og uavhengigheit utgjer eit nullsumspel: tett kopla agentar deler format og snevrar inn det utforskbare rommet; lauskopla agentar vernar autonomien, men ofrar koordinering.}

\prop{6.141}{t}{Systemfasen dikterer det optimale koplingsregimet. Utforsking krev lause koplingar; utnytting krev tette.}

\prop{6.15}{o}{Latens er bandbreiddas temporale dimensjon. Høg latens tvingar agenten til å modellere dei andres framtidige tilstandar. Latensen er ikkje berre ein teknisk eigenskap ved kanalen, men ein epistemisk eigenskap ved agentens situasjon.}

\prop{6.16}{t}{Designaren er aldri åleine i landskapet. Ho navigerer i eit felt av konkurrerande agens frå mikrobielt vev til makroskopiske marknadskrefter. Kompetansen hennar ligg i å koordinere desse uavhengige lyskjeglene til eit stabilt form-optimum.}

\prop{6.2}{o}{Robust kopling emulerer eit overordna system. Overordna system har vidare lyskjegler og samansette mål. Tett samanbinding etablerer globale tilbakekoplingssløyfer — heilskapen vert ein agent. Multi-agent-system oppnår metastabilitet utan lokal optimalisering; denne stabiliteten let tradisjonar persistere gjennom fluktuasjonar som ville ha utsletta delane deira om dei opererte i isolasjon.}

\prop{6.3}{t}{Stilar er deskriptive skuggar. Dei manglar kausal kraft og markerer mellombelse stasjonar i vandringa. Å formgje «i ein stil» er ein kategorifeil — det er imitasjon av statistiske symptom medan årsakstrykka ignorerast. Stilimitasjon kopierer effektar, aldri årsaker. Forma liknar originalen, men sviktar under dei reelle trykka.}

\prop{6.4}{t}{Nyheitsraten avheng av det tilstøytande moglege. Realisering opnar nye regionar og ekspanderer moglegheitsrommet. Formhistoria tømmer ikkje eit lukka repertoar — ho ekspanderer repertoaret gjennom rørsle.}

\prop{6.5}{t}{Inga form er endeleg. Seleksjonstrykk endrar seg; kvar balanse vert med tida suboptimal. Forma er utelukkande gyldig for aktualiseringsaugeblikket. Robustheit krev høg agensstratifisering, aldri rigiditet. Berre den generative strukturen varer: formrom, trykk, landskap, agentar. Formene er flyktige spor. Grammatikken er stabil.}

\prop{6.6}{o}{Rammeverket ytrar ingen ontologisk påstand. Det er eit koordinatsystem for formell spesifikasjon. Substrat-uavhengigheit indikerer fruktbarheit, ikkje absolutt sanning. Agenten er innbakt i formverda han skildrar — det finst inga utside. Formlæra er lokal og situert, aldri transcendent.}

\prop{6.61}{o}{Formgjeving definerer måltilstanden og vel landskapets krefter. Det utvidar lyskjegla; latente affordansar vert synlege. Forma trer fram utan at nokon einskild agent har diktert ho. Vilkåra for kompetent operasjon vert lagt til rette.}

\prop{6.7}{o}{Formverda er det tilfellet som verkar. Navigasjonen held fram.}

\prop{6.71}{o}{Traktaten er ein eigen posisjon i formrommet. Falsifiserbarheit garanterer tekstens gyldigheit.}


\silentchapter{7 Om det som ikkje lèt seg forme}{om det som ikkje lèt seg forme}
\prop{7}{}{Om det som manglar form, lyt ein teia.}
